% Source fingerprint: 2d7726f5ad065ab32b545d9f6c427c768d7b073529961660b05300370c614eaa
% T18: raw TeX cells preserved; no numeric highlighting.
\begin{table}[htbp]
\centering\small\renewcommand{\arraystretch}{1.18}\setlength{\tabcolsep}{4pt}
\caption[流的对象、质量与边界条件]{流的对象、质量与边界条件。这不是无条件的线性包含链。全域与局部版本、实重数与整数重数、弱收敛与平坦收敛分别判断。}
\label{b1:tab:visual-t18}
\begin{tabularx}{\linewidth}{@{}>{\raggedright\arraybackslash}p{3.1cm}>{\raggedright\arraybackslash}p{4.7cm}>{\raggedright\arraybackslash}X@{}}
\toprule
对象 & 质量或几何条件 & 边界与收敛的限制 \\
\midrule
一般流 & 测试微分形式上的连续线性泛函 & 未必具有有限质量；零维流也不必是测度 \\
有限质量流 & \(\mathbf M(T)<\infty\)；可由向量测度表示 & 边界质量仍可能无限 \\
正规流 & \(\mathbf M(T)+\mathbf M(\partial T)<\infty\) & 并不自动可整流或具有整数重数 \\
可整流流（本书约定） & 在可数可整流集上有可测取向与局部可积实重数 & 局部质量有限；全域质量、边界质量与整数重数均需另列 \\
积分流 & \(T\) 与边界都具有整数重数可整流表示，且两者全域质量有限 & 本书把边界可整流性列入定义；局部版本另行说明 \\
有限平坦范数 & 可由 \(T=A+\partial B\) 的有限质量分解控制 & 不自动有有限质量；薄矩形边界可平坦趋零而质量不趋零 \\
\bottomrule
\end{tabularx}
\end{table}
